\section{Matched flattened coordinates}
\label{sec:app_ft_matched_coordinates}

Let $\beta:J(Q)\simeq J(P)$ match the basis sectors and let $\tau_k$ match the
copies in sector $k$. Flatten the pair $(k,q)$ to the copy index of $Q$. The
matched $P$-data on that index are the weight and tensor
\begin{align}
    \widetilde\mu_{k,q}
     & =\mu_{\beta(k),\tau_k(q)},
    &
    \widetilde P_{k,q}
     & =P_{\beta(k)},
    \notag
\end{align}
where the second expression is read after the matched bond-dimension
identification. The block gauge repeats $X_k$ on every copy $q$.

\begin{definition}[Matched weights, tensors, and repeated gauges]
    \label{def:ft_matched_flat_data}
    \lean{MPSTensor.matched_p_weight,
        MPSTensor.matched_p_basis,
        MPSTensor.matched_block_gauge}
    \leanok
    The matched flattened data are the functions
    $\widetilde\mu$, $\widetilde P$, and $(k,q)\mapsto X_k$ defined above on
    the flattened copy coordinates of $Q$.
\end{definition}

\begin{lemma}[Equality of the matched total bond dimensions]
    \label{lem:ft_total_dim_eq_of_match}
    \lean{MPSTensor.SectorDecomposition.totalDim_eq_of_match}
    \leanok
    \uses{def:ft_matched_flat_data}
    If matched sectors have equal bond dimensions and matched copy sets are
    bijective, then $P$ and $Q$ have the same total bond dimension.
\end{lemma}

\begin{proof}\leanok
    Reindex first by the sector bijection and then by each copy bijection.
    Every summand $\dim P_{\beta(k)}$ becomes $\dim Q_k$, so
    \begin{align}
        \sum_{j\in J(P)}r_j(P)\dim P_j
         & =\sum_{k\in J(Q)}r_k(Q)\dim Q_k.
        \notag
    \end{align}
\end{proof}

\begin{definition}[Sector and bond-coordinate equivalences]
    \label{def:ft_sector_flat_equiv}
    \lean{MPSTensor.SectorDecomposition.sectorFlatSigmaEquiv,
        MPSTensor.SectorDecomposition.sectorFlatDimEquiv}
    \leanok
    \uses{lem:ft_total_dim_eq_of_match}
    The sector bijection, copy bijections, and bond-dimension equalities induce
    an equivalence of the dependent triples $(j,q,a)$ and $(k,q',a')$. After
    flattening the dependent bond coordinate, this gives an equivalence of the
    standard total bond-index sets.
\end{definition}

\begin{proof}\leanok
    Send $(k,q,a)$ to
    $(\beta(k),\tau_k(q),a)$, using the bond-dimension equality to identify the
    last coordinate. The inverse uses $\beta^{-1}$ and $\tau_k^{-1}$. Both
    composites are the identity componentwise. Conjugating by the standard
    equivalence between dependent triples and a finite interval gives the
    flattened bond-coordinate equivalence.
\end{proof}

\begin{theorem}[Equal-MPV gauge with all matched witnesses]
    \label{thm:sector_bnt_equal_mps_gaugeEquiv_witnesses}
    \lean{MPSTensor.ft_sector_bnt_equal_mps_gaugeEquiv_witnesses}
    \leanok
    \uses{thm:sector_bnt_equal_global_gauge,
        lem:ft_total_dim_eq_of_match, def:ft_matched_flat_data}
    Under the hypotheses of the equal-MPV theorem, one may choose
    simultaneously the sector bijection, bond-dimension equalities, copy
    bijections, phases, block gauges, total-dimension equality, and the
    flattened global gauge. They satisfy the matched block and weight
    identities of Theorem~\ref{thm:sector_bnt_equal_global_gauge}, together
    with the conjugation of $Q$ by the matched-coordinate $P$-tensor.
\end{theorem}

\begin{proof}\leanok
    Theorem~\ref{thm:sector_bnt_equal_global_gauge} supplies all witnesses
    except the displayed equality of total dimensions. Apply
    Lemma~\ref{lem:ft_total_dim_eq_of_match} to its sector and copy bijections.
    The resulting equality identifies the total bond spaces, while the global
    conjugation equation is unchanged.
\end{proof}

\section{Permutation gauges and literal coordinates}
\label{sec:app_ft_permutation_gauges}

The matched $P$-tensor has the correct blocks but lists them in $Q$-order. A
permutation matrix restores the literal order of $P$. We record its invertible
and unitary forms before proving the coordinate conversion used in
Theorem~\ref{thm:sector_bnt_equal_mps_gaugeEquiv_literal}.

\begin{definition}[Invertible permutation gauge]
    \label{def:ft_perm_gl}
    \lean{MPSTensor.permGL}
    \leanok
    For a permutation $\rho$ of a finite bond-index set, let $\Pi_\rho$ be its
    permutation matrix. The corresponding element of the general linear group
    is denoted $G_\rho$.
\end{definition}

\begin{lemma}[Matrix, inverse, and unitary identities for a permutation gauge]
    \label{lem:ft_perm_gl_identities}
    \lean{MPSTensor.permGL_val,
        MPSTensor.permGL_inv_val,
        MPSTensor.permGL_mem_unitaryGroup}
    \leanok
    \uses{def:ft_perm_gl}
    The matrix of $G_\rho$ is $\Pi_\rho$, its inverse is
    $\Pi_{\rho^{-1}}$, and it is unitary.
\end{lemma}

\begin{proof}\leanok
    Every row and every column of $\Pi_\rho$ contains exactly one nonzero
    entry, equal to $1$. Hence
    \begin{align}
        \Pi_\rho\Pi_{\rho^{-1}}
         & =\Id,
        &
        \Pi_\rho^\dagger
         & =\Pi_{\rho^{-1}}.
        \notag
    \end{align}
    These identities give both the inverse formula and unitarity.
\end{proof}

The coordinate equivalence of Definition~\ref{def:ft_sector_flat_equiv}
induces a permutation $\rho$ after the total-dimension identification. For any
matrix $M$, conjugation by its permutation gauge reindexes both matrix
coordinates:
\begin{align}
    G_\rho M G_\rho^{-1}
     & =M\circ(\rho^{-1}\times\rho^{-1}).
    \label{eq:app_ft_perm_conjugation}
\end{align}
Applying (\ref{eq:app_ft_perm_conjugation}) to $P_{\mathrm{tot}}^i$ gives the
matched-coordinate tensor. Composing $G_\rho$ with the block-diagonal gauge
therefore gives the literal gauge of
Theorem~\ref{thm:sector_bnt_equal_mps_gaugeEquiv_literal}. This also explains
why the total-dimension cast and the sector permutation must occur in that
order.

\section{Unitary witness refinements}
\label{sec:app_ft_unitary_witnesses}

The main text proves that a specified gauge between left-canonical irreducible
sectors can be replaced by a unitary without changing its phase. We record the
existential form and then repeat the construction over every matched copy.

\begin{lemma}[Unitary gauge for left-canonical irreducible tensors]
    \label{lem:left_canonical_irreducible_unitary_gauge}
    \lean{MPSTensor.exists_unitaryConj_gaugePhase_of_leftCanonical_irreducible}
    \leanok
    \uses{def:left_canonical_tensor, def:gauge_phase_equiv,
        def:irreducible_tensor,
        lem:left_canonical_irreducible_same_phase_unitary_gauge}
    Let $A$ and $B$ be left-canonical irreducible tensors of common positive bond
    dimension that are gauge-phase equivalent. Then the gauge can be taken
    unitary: there are a unitary matrix $U$ and a scalar $\zeta$ with
    $|\zeta|=1$ such that, for every physical index~$i$,
    \begin{align}
        B^i&=\zeta\,UA^iU^\dagger. \notag
    \end{align}
\end{lemma}

\begin{proof}\leanok
    \uses{lem:left_canonical_irreducible_same_phase_unitary_gauge}
    Unpack the gauge-phase equivalence as
    $B^i=\zeta XA^iX^{-1}$ with $X$ invertible and $\zeta\ne0$.
    Lemma~\ref{lem:left_canonical_irreducible_same_phase_unitary_gauge}
    replaces $X$ by a unitary $U$ while retaining this same phase and proves
    $|\zeta|=1$.
\end{proof}

\begin{lemma}[Unitary gauge in matched sector coordinates]
    \label{lem:sector_bnt_equal_unitary_matched_gauge}
    \lean{MPSTensor.ft_sector_bnt_equal_unitary_global_gauge_witnesses}
    \leanok
    \uses{lem:sector_bnt_equal_matched_copy_weight_witnesses,
        lem:left_canonical_irreducible_same_phase_unitary_gauge,
        thm:sector_bnt_global_gauge_of_copy_weight_matching,
        lem:unitary_global_gauge_of_blocks}
    Let $P$ and $Q$ be BNT canonical forms whose total tensors generate the
    same MPV family at every positive length. There are a sector bijection
    $\beta$, copy bijections $\tau_k$, phases $\zeta_k$, sector unitaries $U_k$,
    and a unitary block-diagonal matrix $X$ such that
    \begin{align}
        Q_k^i
        &= \zeta_k U_kP_{\beta(k)}^iU_k^\dagger,
        \notag\\
        \nu_{k,q}
        &= \zeta_k^{-1}\mu_{\beta(k),\tau_k(q)},
        \notag\\
        |\zeta_k|
        &= 1.
        \notag
    \end{align}
    The matched sectors and copies have equal bond dimensions and
    multiplicities. If $P_{\mathrm{match}}$ denotes the corresponding
    reordered block assembly, then
    $Q_{\mathrm{flat}}^i=XP_{\mathrm{match}}^iX^{-1}$.
\end{lemma}
\begin{proof}\leanok
    \uses{lem:sector_bnt_equal_matched_copy_weight_witnesses,
        lem:left_canonical_irreducible_same_phase_unitary_gauge,
        lem:unitary_global_gauge_of_blocks}
    Apply Lemma~\ref{lem:left_canonical_irreducible_same_phase_unitary_gauge}
    to each matched sector, retaining the phase in the copy-weight identity.
    Repeating $U_k$ on every copy gives $X=\bigoplus_{k,q}U_k$, which is unitary
    by Lemma~\ref{lem:unitary_global_gauge_of_blocks}. For every matched copy
    $(k,q)$,
    \begin{align}
        \nu_{k,q}Q_k^i
        &= \zeta_k^{-1}\mu_{\beta(k),\tau_k(q)}
          \zeta_k U_kP_{\beta(k)}^iU_k^\dagger
        = \mu_{\beta(k),\tau_k(q)}U_kP_{\beta(k)}^iU_k^\dagger.
        \notag
    \end{align}
    These are precisely the blocks of
    $Q_{\mathrm{flat}}^i=XP_{\mathrm{match}}^iX^{-1}$.
\end{proof}
